\section*{Capítulo \ref{ch:g11:igneous-rocks} --- Las rocas ígneas}

\begin{solution}{exo:g11:igneous-rocks:1}
Una \omterm{def:g11:igneous-rocks:volcanic-plutonic}{roca volcánica} se formó a partir de un magma que se solidificó en la
superficie, como lava; una \omterm{def:g11:igneous-rocks:volcanic-plutonic}{roca plutónica} cristalizó en profundidad, dentro de
otras rocas. Parejas: basalto y \omterm{def:g11:ocean-ridges:gabbro}{gabro} (\omterm{def:g11:igneous-rocks:felsic-mafic}{máficas}), \omterm{def:g11:igneous-rocks:rock-names}{andesita} y \omterm{def:g11:igneous-rocks:rock-names}{diorita}
(\omterm{def:g11:igneous-rocks:felsic-mafic}{intermedias}), \omterm{def:g11:igneous-rocks:rock-names}{riolita} y granito (\omterm{def:g11:igneous-rocks:felsic-mafic}{félsicas}).
\end{solution}

\begin{solution}{exo:g11:igneous-rocks:2}
(1) Vítrea. (2) De grano grueso. (3) \omterm{def:g11:igneous-rocks:porphyritic}{Porfídica} (\omterm{def:g11:igneous-rocks:porphyritic}{fenocristales} de olivino en
una matriz fina). (4) De grano fino.
\end{solution}

\begin{solution}{exo:g11:igneous-rocks:3}
\qty{48}{\percent}: \omterm{def:g11:igneous-rocks:felsic-mafic}{máfico}, \omterm{def:g11:ocean-ridges:gabbro}{gabro}. \qty{60}{\percent}: \omterm{def:g11:igneous-rocks:felsic-mafic}{intermedio}, \omterm{def:g11:igneous-rocks:rock-names}{diorita}.
\qty{72}{\percent}: \omterm{def:g11:igneous-rocks:felsic-mafic}{félsico}, granito.
\end{solution}

\begin{solution}{exo:g11:igneous-rocks:4}
Un \omterm{def:g11:igneous-rocks:intrusion}{dique} es una lámina de roca ígnea que corta las capas de las rocas en las
que se intruye; un \omterm{def:g11:igneous-rocks:intrusion}{filón capa} es una lámina paralela a las capas, situada
entre ellas; un \omterm{def:g11:igneous-rocks:intrusion}{plutón} es una \omterm{def:g11:igneous-rocks:intrusion}{intrusión} grande, de cualquier forma, formada en
profundidad. El \omterm{def:g11:igneous-rocks:intrusion}{dique} vertical que sube desde el \omterm{def:g11:igneous-rocks:intrusion}{plutón} alimentaba el volcán.
\end{solution}

\begin{solution}{exo:g11:igneous-rocks:5}
El olivino se forma primero; el feldespato potásico, el cuarzo y a veces la
moscovita, al final. La plagioclasa empieza siendo rica en calcio y se hace
cada vez más rica en sodio a medida que el magma se enfría.
\end{solution}

\begin{solution}{exo:g11:igneous-rocks:6}
(1) A: un \qty{74}{\percent} de \ce{SiO2}, \omterm{def:g11:igneous-rocks:felsic-mafic}{félsica}; B: un \qty{47}{\percent},
\omterm{def:g11:igneous-rocks:felsic-mafic}{máfica}. (2) A es pobre en \ce{FeO}, \ce{MgO} y \ce{CaO} y rica en \ce{Na2O} y
\ce{K2O}, como un magma \omterm{def:g11:igneous-rocks:felsic-mafic}{félsico}; B es rica en \ce{FeO}, \ce{MgO} y \ce{CaO}, y
pobre en \ce{Na2O}, como un magma \omterm{def:g11:igneous-rocks:felsic-mafic}{máfico}. (3) A es una \omterm{def:g11:igneous-rocks:rock-names}{riolita}; B, un \omterm{def:g11:ocean-ridges:gabbro}{gabro}.
\end{solution}

\begin{solution}{exo:g11:igneous-rocks:7}
El tiempo crece como el cuadrado del espesor. \qty{1}{m}: $10/4 = 2.5$
días. \qty{20}{m}: $10 \times 10^2 = 1000$ días, unos 2.7 años.
\qty{200}{m}: $10 \times 100^2 = \num{100000}$ días, unos 270 años.
\end{solution}

\begin{solution}{exo:g11:igneous-rocks:8}
(1) \unit{m^2} dividido entre \unit{m^2/s} da segundos. (2) $400/10^{-6} =
\qty{4e8}{s}$, es decir, $4 \times 10^8 / 3.16 \times 10^7 \approx 12.7$
años. (3) $0.21 \times 12.7 \approx 2.7$ años.
\end{solution}

\begin{solution}{exo:g11:igneous-rocks:9}
Es \omterm{def:g11:igneous-rocks:porphyritic}{porfídica}: \omterm{def:g11:igneous-rocks:porphyritic}{fenocristales} de plagioclasa en una matriz de grano fino.
Primero el magma se enfrió despacio en profundidad, y los cristales de
plagioclasa crecieron mucho mientras la mayor parte del magma seguía líquida;
luego ascendió y se enfrió deprisa, y el líquido restante cristalizó formando
la matriz fina.
\end{solution}

\begin{solution}{exo:g11:igneous-rocks:10}
(1) A la izquierda: 2 Mg, 2 Si, 6 O; a la derecha: 2 Mg, 2 Si, 6 O.
(2) \qty{1}{kg} de olivino son $1000/140.7 \approx 7.11$~mol, que toman otros
tantos moles de sílice: $7.11 \times 60.1 \approx \qty{427}{g}$. (3) Si no
queda sílice en el líquido, la reacción no puede continuar: el olivino se
queda.
\end{solution}

\begin{solution}{exo:g11:igneous-rocks:11}
(1) $2.5 \times 20 \times 1100 = \qty{55000}{MJ}$, \qty{5.5e10}{J}.
(2) $900 \times 20 = \qty{18000}{MJ}$, \qty{1.8e10}{J}. (3) El calor de
cristalización es alrededor de un tercio del calor perdido al enfriarse, y una
cuarta parte del total: basta para frenar mucho la solidificación, porque se
libera todo por encima del solidus.
\end{solution}

\begin{solution}{exo:g11:igneous-rocks:12}
(1) Los cristales más pequeños, en el contacto, que se solidificó en menos de
un día; los más grandes, en el centro, que siguió líquido unos 950 días.
(2) No: los primeros \qty{5}{m} tardaron unos 545 días, y los últimos
\qty{5}{m}, unos 400 días; el frente se frena al principio y luego se acelera
cerca del centro, donde se encuentra el calor que sale por los dos lados.
(3) Incluso en su centro, el \omterm{def:g11:igneous-rocks:intrusion}{filón capa} se solidificó en menos de tres años;
un granito de cristales grandes necesita mucho más tiempo.
\end{solution}

\begin{solution}{exo:g11:igneous-rocks:13}
(1) El \omterm{def:g11:igneous-rocks:intrusion}{plutón} es 100 veces más grueso: $2.6 \times 100^2 = \num{26000}$
años. (2) Por ejemplo: un \omterm{def:g11:igneous-rocks:intrusion}{plutón} profundo está en una roca que ya está
caliente, así que pierde calor más despacio que hacia una roca a
\qty{50}{\degreeCelsius}; puede recibir magma nuevo una y otra vez; un \omterm{def:g11:igneous-rocks:intrusion}{plutón}
no es una lámina infinita; las propiedades de la roca cambian con la
temperatura.
\end{solution}

\begin{solution}{exo:g11:igneous-rocks:14}
Un magma \omterm{def:g11:igneous-rocks:felsic-mafic}{máfico} está caliente y es pobre en sílice: cristaliza los primeros
minerales de la serie ---olivino, piroxeno, plagioclasa rica en calcio--- y se
queda sin sílice antes de que pueda formarse cuarzo. Un magma \omterm{def:g11:igneous-rocks:felsic-mafic}{félsico} está
más frío y es rico en sílice: en él el olivino no es estable, y llega hasta el
final de la serie, formando plagioclasa rica en sodio, feldespato potásico y
cuarzo.
\end{solution}

\begin{solution}{exo:g11:igneous-rocks:15}
(1) $b = \sqrt{t\kappa/C} = \sqrt{86400 \times 10^{-6}/0.21} \approx
\qty{0.64}{m}$. (2) La ley da $0.21 \times 25/10^{-6} \approx
\qty{5.3e6}{s}$, unos dos meses: del mismo orden que un mes (una colada de
lava también pierde calor hacia el aire por su parte superior). (3) Vítrea en
el material lanzado al aire o al agua; de grano fino en las coladas de lava y
las \omterm{def:g11:minerals:thin-section}{láminas delgadas}; más gruesa, pero todavía fina, en las intrusiones
gruesas y poco profundas.
\end{solution}

\begin{solution}{pb:g11:igneous-rocks:1}
\textbf{1.} Un \omterm{def:g11:igneous-rocks:intrusion}{filón capa}: es paralelo a los estratos y está entre ellos.

\textbf{2.} En el contacto, de grano muy fino, casi vítrea; en el centro, de
grano fino a medio, con los minerales apenas visibles con una lupa.

\textbf{3.} El contacto perdió su calor más deprisa, hacia la arenisca fría:
los cristales no tuvieron tiempo de crecer.

\textbf{4.} Un borde enfriado. Demuestra que el magma estaba caliente cuando
se emplazó, y la arenisca, fría.

\textbf{5.} Un \omterm{def:g11:igneous-rocks:intrusion}{dique}, el conducto de alimentación del \omterm{def:g11:igneous-rocks:intrusion}{filón capa}.

\textbf{6.} \omterm{def:g11:igneous-rocks:felsic-mafic}{Máfico}: un \qty{49}{\percent} de \ce{SiO2}, rico en hierro,
magnesio y calcio.

\textbf{7.} Olivino, piroxeno y plagioclasa rica en calcio.

\textbf{8.} Un basalto; la \omterm{def:g11:igneous-rocks:volcanic-plutonic}{roca plutónica} de la misma composición es el \omterm{def:g11:ocean-ridges:gabbro}{gabro}.

\textbf{9.} La \omterm{def:g11:igneous-rocks:texture}{textura} depende del ritmo de enfriamiento, que es distinto del
contacto al centro, no de la composición.

\textbf{10.} Primero el olivino; la plagioclasa empieza más o menos en el
mismo punto que el piroxeno.

\textbf{11.} $b^2/\kappa = 400/10^{-6} = \qty{4e8}{s}$.

\textbf{12.} $0.206 \times 4 \times 10^8 \approx \qty{8.2e7}{s}$, es decir,
$8.2 \times 10^7 / 3.16 \times 10^7 \approx 2.6$ años.

\textbf{13.} $0.206 \times 4/10^{-6} \approx \qty{8.2e5}{s}$, unos 9.5
días.

\textbf{14.} $0.07 \times 4 \times 10^8 = \qty{2.8e7}{s}$, unos 0.9 años. El
magma que cristaliza libera calor, que también hay que evacuar antes de que el
magma se vuelva sólido.

\textbf{15.} Unos 950 días frente a 0.6 días: unas 1600 veces más deprisa.

\textbf{16.} Unos 950 días, unos 2.6 años: mucho más que una colada de
lava, dentro de los meses o años de las intrusiones poco profundas.

\textbf{17.} Mil veces más tiempo exige un espesor $\sqrt{1000} \approx
32$
veces mayor: unos \qty{630}{m}.

\textbf{18.} Cerca de la superficie, un cuerpo de magma saldría en erupción, o
se enfriaría demasiado deprisa; un cuerpo grande que se enfría durante muchísimo
tiempo tiene que estar en profundidad, en una roca que ya está caliente.

\textbf{19.} Las ha desgastado la erosión.

\textbf{20.} El \omterm{def:g11:igneous-rocks:intrusion}{filón capa} tardó unos \textbf{2.6 años} en solidificarse.
\end{solution}
